%%%%%%%%%%%%%%%%%%%%%%%%%%%%%%%%%%%%%%%%%%%%%%%%%%%%%%%%%%%%%%%%%%%%%%
% Lớp 11 | Đường thẳng, mặt phẳng. Quan hệ song song trong | Hình lăng trụ và hình hộp (xiên)
% Dạng 1H4?5-4: Toán thực tế áp dụng hình lăng trụ và hình hộp
% Số câu: 5  (sinh tự động bởi tach_cau.py; từ nay thầy sửa trực tiếp file này)
%%%%%%%%%%%%%%%%%%%%%%%%%%%%%%%%%%%%%%%%%%%%%%%%%%%%%%%%%%%%%%%%%%%%%%

%%%%%%%%%%%%%%%% TRẮC NGHIỆM 4 PHƯƠNG ÁN %%%%%%%%%%%%%%%%

% nguồn: Ok_TracNghiemLop11_Lan1
\begin{ex}%[1H4K5-4]%[ma:c5301093]
	Cho tứ diện $ABCD$ đều có  cạnh bằng $2\sqrt{2}$. Gọi $G$ là trọng tâm tứ diện $ABCD$ và $M$ là trung điểm $AB$. Khoảng cách giữa hai đường thẳng $BG$ và $CM$ bằng
	\choice 
		{\True  $\dfrac{2}{\sqrt{14}}$}
		{$\dfrac{2}{\sqrt{5}}$} 
		 {$\dfrac{3}{2\sqrt{5}}$}
		 {$\dfrac{2}{\sqrt{10}}$}
	\loigiai{ 
		\begin{center}
			\begin{tikzpicture}[scale=1,font=\footnotesize,line join=round,line cap=round,>=stealth]
					 \def\b{7}
					 \def\c{5}
					\def\gA{-22} 
					 \def\h{4}
					\path (0,0)  coordinate (B)
						 (0:\b) coordinate (D)
						(\gA:\c)  coordinate (C)
						($(D)!0.5!(C)$)  coordinate (N)
						 ($(C)!0.5!(N)$) coordinate (K)
						($(B)!2/3!(N)$) coordinate  (H)
						 ($(H)+(78:\h)$) coordinate (A)
						($(A)!0.5!(B)$) coordinate  (M)
						 ($(M)!0.5!(N)$) coordinate (G)
						 ($(B)!0.55!(K)$) coordinate (I)
						($(I)!(H)!(G)$)  coordinate (J);
					\foreach \i in {D,B,C}{\draw  (A)--(\i);}
					\draw  (B)--(C)--(D)
						 (M)--(C);
					 \draw[dashed] (D)--(B)--(N)
						(G)--(B)--(K)--cycle  (A)--(H)
						(M)--(N) 
						 (J)--(H)--(I)--(G);
					\pic[draw,angle radius=2mm] {right angle  = B--N--C};
					\pic[draw,angle radius=2mm] {right angle =  I--J--H};
					\pic[draw,angle  radius=2mm] {right angle = K--I--H};
					\foreach \i/\g in  {A/90,B/-135,C/-45,D/-45,H/45,K/-45,I/-135,N/-45,M/135,G/40,J/135}{\draw[fill](\i) circle (1pt)
						 ($(\i)+(\g:3mm)$) node[scale=1]{$\i$};}
				\end{tikzpicture}
		\end{center}
		Gọi $N$ là trung điểm $CD$, khi đó $G$ là trung điểm $MN$ và $AG$ đi qua trọng  tâm $H$ của tam giác $BCD$.\\
		Ta có $AH\perp(BCD$.  và $AH=\sqrt{AB^2-BH^2}=\sqrt{(2\sqrt{2})^2-\left(\dfrac{2\sqrt{6}}{3}\right)^2}=\dfrac{4\sqrt{3}}{3}$.\\
		 Ta có $GH=\dfrac{1}{4} AH=\dfrac{\sqrt{3}}{3}$.\\
		Gọi $K$ là trung điểm $CN$ thì $GK\parallel CM$ nên  $CM\parallel(BGK)$.\\
		 Do đó
		$\mathrm{d}(BG; CM)=\mathrm{d}(C;(BGK))=\mathrm{d}(N;(BGK))=\dfrac{3}{2}  \mathrm{d}(H;(BGK))$.\\
		Kẻ $HI\perp BK$,  $HJ\perp GI$ với $I\in BK$, $J\in GI$.\\
		Khi  đó $HJ\perp(BGK)$ và $HJ=\mathrm{d}(H;(BGK))$.\\
		Ta  có $BK=\sqrt{BN^2+NK^2}=\sqrt{(\sqrt{6})^2+\left(\dfrac{\sqrt{2}}{2}\right)^2}=\dfrac{\sqrt{26}}{2}$.\\
		Ta có  $HI=BH\cdot \sin \widehat{KBN}=BH\cdot \dfrac{KN}{BK}=\dfrac{2\sqrt{6}}{3} \cdot \dfrac{\dfrac{\sqrt{2}}{2}}{\dfrac{\sqrt{26}}{2}}=\dfrac{2\sqrt{6}}{3\sqrt{13}}$.\\
		 Do đó $HJ=\dfrac{HI\cdot HG}{\sqrt{HI^2+HG^2}}=\dfrac{\dfrac{2\sqrt{6}}{3\sqrt{13}} \cdot \dfrac{\sqrt{3}}{3}}{\sqrt{\left(\dfrac{2\sqrt{6}}{3\sqrt{13}}\right)^2+\left(\dfrac{\sqrt{3}}{3}\right)^2}}=\dfrac{2\sqrt{2}}{3\sqrt{7}}$.\\
		Vậy $\mathrm{d}(BG; CM)=\dfrac{3}{2} \mathrm{d}(H;(BGK))=\dfrac{3}{2} HJ=\dfrac{3}{2} \cdot  \dfrac{2\sqrt{2}}{3\sqrt{7}}=\dfrac{2}{\sqrt{14}}$.
	 }
\end{ex}

% nguồn: Ok_TracNghiemLop11_Lan1
\begin{ex}%[1H4B5-4]%[ma:59222c9a]
	Cho hình chóp $S.ABCD$ có đáy là hình chữ nhật với $AB=a$, $AD=2a$. Biết $SA=2a$ và $SA$ vuông góc với mặt phẳng  đáy. Khoảng cách giữa hai đường thẳng $BD$ và $SC$ bằng
	\choice 
		{$\dfrac{a \sqrt{30}}{6}$} 
		{\True $\dfrac{2\sqrt{21} a}{21}$} 
		{$\dfrac{4\sqrt{21} a}{21}$} 
		 {$\dfrac{a \sqrt{30}}{12}$}
	 \loigiai{
		\begin{center}
			\begin{tikzpicture}[scale=1,font=\footnotesize,line join=round,line cap=round,>=stealth]
					 \def\a{4}
					 \def\b{2}
					 \def\g{35}
					\def\h{3.5} 
					 \path (0,0) coordinate (B)
						(0:\a) coordinate  (C)
						(\g:\b) coordinate (A) 
						(A)+(0:\a)  coordinate (D)
						 ($(B)!0.5!(D)$) coordinate (O)
						(A)+(90:\h)  coordinate (S)
						 ($(S)!0.5!(A)$) coordinate (M);
					\foreach \i in  {B,C,D}{\draw (S)--(\i);}
					\draw  (B)--(C)--(D);
					\draw[dashed] (S)--(A)--(B)--(D)--(A)--(C) 
						 (B)--(M)--(D)
						(O)--(M); 
					 \pic[draw,angle radius=2mm] {right angle = C--A--S};
					\foreach  \i/\g in {S/90,A/135,B/-135,C/-45,D/45,O/-90,M/170}{\draw[fill](\i) circle (1pt)
						 ($(\i)+(\g:3mm)$) node[scale=1]{$\i$};}
				\end{tikzpicture}
		\end{center}
		Gọi $O$ là tâm  hình chữ nhật và $M$ là trung điểm $SA$, ta có $SC\parallel(BMD)$.\\
		Do đó  $\mathrm{d}(SC, BD)=\mathrm{d}(SC, (BMD))=\mathrm{d}(S, (BMD))=\mathrm{d}(A, (BMD))=h$.\\
		Ta có $AM$, $AB$, $AD$ đôi  một vuông góc nên
		$$ 
		\dfrac{1}{h^2}=\dfrac{1}{AM^2}+\dfrac{1}{AB^2}+\dfrac{1}{AD^2}=\dfrac{4}{a^2}+\dfrac{1}{a^2}+\dfrac{1}{4a^2}. 
		 $$
		Vậy  $\mathrm{d}(SC, BD)=\dfrac{2a \sqrt{21}}{21}$.
	 }
\end{ex}


%%%%%%%%%%%%%%%% ĐÚNG/SAI %%%%%%%%%%%%%%%%

% nguồn: Ok_dungsai_lop11_Lan1
\begin{ex}%[1H4B5-4]%[ma:64176ac3]
 Cho hình chóp đều $S.ABC$ có cạnh đáy bằng $a$, gọi $O$  là tâm của đáy và $SO=\dfrac{a \sqrt{3}}{3}$. Khi đó
 \choiceTF 
   {$AO=\dfrac{a\sqrt{3}}{2}$}
  {\True  $\mathrm{d}(O, SA)=\dfrac{a\sqrt{6}}{6}$}
  {\True  Kẻ đường cao $AI$ của tam giác $ABC$, khi đó $OI=\dfrac{a\sqrt{3}}{6}$}
  {$\mathrm{d}(O,(SBC))=\dfrac{a\sqrt{15}}{12}$} 
 \loigiai{ 
  \begin{center}
   \begin{tikzpicture}[scale=1,font=\footnotesize,line join=round,line cap=round,>=stealth]
      \def\b{4.5}
      \def\c{3.5}
      \def\gA{-30}
     \def\h{3.5} 
     \path (0,0) coordinate  (A)
      (0:\b) coordinate (C) 
      (\gA:\c) coordinate  (B)
      ($(B)!0.5!(C)$) coordinate (I) 
       ($(A)!2/3!(I)$) coordinate (O)
      (O)+(90:\h) coordinate (S) 
       ($(S)!(O)!(A)$) coordinate (H)
      ($(S)!(O)!(I)$) coordinate  (K);
     \foreach  \i in {A,B,C}{\draw (S)--(\i);}
     \draw (A)--(B)--(C) 
       (S)--(I);
     \draw[dashed]  (A)--(C)
       (A)--(I)
       (H)--(O)--(K)
      (S)--(O); 
     \pic[draw,thin,angle radius=2mm] {right angle  = O--H--S};
     \pic[draw,thin,angle radius=2mm] {right angle = S--O--A}; 
     \pic[draw,thin,angle radius=2mm] {right  angle = S--K--O};
     \pic[draw,thin,angle  radius=2mm] {right angle = A--I--B};
     \foreach \i/\g in  {S/90,A/-135,B/-90,C/-45,O/-90,H/150,K/0,I/0}{\draw[fill=black](\i) circle (1pt)
       ($(\i)+(\g:3mm)$) node[scale=1]{$\i$};}
    \end{tikzpicture}
  \end{center}
  \begin{itemchoice}
   \itemch Sai. Kẻ đường cao $AI$ của  tam giác $ABC$, ta có $O$ thuộc $AI$.\\
   Trong mặt phẳng $(SAI)$, dựng $OH\perp SA$ tại  $H\Rightarrow \mathrm{d}(O, SA)=OH$.\\
   Tam giác $ABC$ đều cạnh $a$  nên $AO=\dfrac{2}{3} AI=\dfrac{2}{3} \cdot \dfrac{a \sqrt{3}}{2}=\dfrac{a \sqrt{3}}{3}=SO$.
   \itemch Đúng. Tam giác $SAO$ vuông cân tại  $O$ nên $OH=\dfrac{SA}{2}=\dfrac{\dfrac{a \sqrt{3}}{3} \cdot \sqrt{2}}{2}=\dfrac{a \sqrt{6}}{6}$.\\
    Vậy $\mathrm{d}(O, SA)=OH=\dfrac{a\sqrt{6}}{6}$.
   \itemch Đúng. Ta xét khoảng cách  từ $O$ đến mặt bên $(SBC)$.\\
   Kẻ đường cao $OK$ của tam giác $SOI\Rightarrow OK\perp SI$. \hfill(1)\\ 
   Ta có $\heva{&BC\perp SO\ (\text{do  } SO\perp(ABC) \\ &BC\perp AI} \Rightarrow BC\perp(SAI) \Rightarrow BC\perp OK$. \hfill(2)\\
   Từ $(1)$ và $(2)$ suy ra $OK\perp(SBC)$  hay $OK=\mathrm{d}(O, (SBC))$.\\
   Ta có ${OI=\dfrac{AI}{3}=\dfrac{\dfrac{a \sqrt{3}}{2}}{3}=\dfrac{a \sqrt{3}}{6}}$. 
   \itemch Sai. Tam giác $SOI$  vuông tại $O$ có đường cao $OK$ nên
   $$\dfrac{1}{OK^2}=\dfrac{1}{SO^2}+\dfrac{1}{OI^2} \Rightarrow OK=\dfrac{SO\cdot OI}{\sqrt{SO^2+OI^2}}=\dfrac{\dfrac{a \sqrt{3}}{3} \cdot \dfrac{a  \sqrt{3}}{6}}{\sqrt{\dfrac{3a^2}{9}+\dfrac{3a^2}{36}}}=\dfrac{a \sqrt{15}}{15}.$$
   Vậy ${\mathrm{d}(O, (SBC.)=OK=\dfrac{a  \sqrt{15}}{15}}$.
  \end{itemchoice}
 } 
\end{ex}

% nguồn: Ok_dungsai_lop11_Lan1
\begin{ex}%[1H4K5-4]%[ma:8d5fb2ee]
 Cho hình chóp $S.ABCD$ có đáy $ABCD$ là hình chữ  nhật với $AB=2a$, $AD=a$. Hình chiếu của $S$ lên mặt phẳng $(ABCD)$ là trung điểm $H$ của $AB$ và $\widehat{SCH}=45^{\circ}$. Khi đó
 \choiceTF 
  {\True  $BC\perp (SAB)$}
  {\True  $\mathrm{d}(H,(SBC))=\dfrac{a\sqrt{6}}{3}$}
  {\True Gọi $K$ là trung điểm $CD$ khi đó $CD\perp (SHK)$} 
  {$\mathrm{d}(H, (SCD))=\dfrac{a\sqrt{6}}{2}$} 
 \loigiai{ 
  \begin{center}
   \begin{tikzpicture}[scale=1,font=\footnotesize,line join=round,line cap=round,>=stealth]
      \def\a{4}
     \def\b{2.5} 
      \def\g{55}
     \def\h{3.5} 
      \path (0,0) coordinate (B)
       (0:\a) coordinate (C)
      (\g:\b)  coordinate (A)
      (A)+(0:\a) coordinate (D) 
      ($(A)!0.5!(B)$) coordinate (H) 
      ($(C)!0.5!(D)$) coordinate  (K)
      (H)+(90:\h) coordinate  (S)
      ($(B)!0.3!(S)$) coordinate (E) 
       ($(S)!(H)!(K)$) coordinate (I);
     \foreach \i in {B,C,D}{\draw  (S)--(\i);}
     \draw  (B)--(C)--(D)
      (S)--(K); 
      \draw[dashed] (S)--(A)--(B)
      (D)--(A) 
       (S)--(H)--(E)
      (C)--(H)--(I) 
      (H)--(K); 
     \pic[draw,thin,angle radius=2mm] {right angle =  A--H--S};
     \pic[draw,thin,angle radius=2mm] {right angle  = H--K--C};
     \pic[draw,thin,angle radius=2mm] {right  angle = H--I--K};
     \pic[draw,thin,angle radius=2mm] {right angle =  B--E--H};
      \pic[draw,angle radius=15pt,fill=gray!20]{angle= S--C--H};
     \foreach \i/\g in  {S/90,A/135,B/-135,C/-45,D/45,E/160,K/-45,I/45,H/-80}{\draw[fill=black](\i) circle (1pt)
      ($(\i)+(\g:3mm)$) node[scale=1]{$\i$};} 
    \end{tikzpicture}
  \end{center}
  \begin{itemchoice}
   \itemch  Đúng. Kẻ đường cao $HE$ trong tam giác $SBH$. \hfill(1)\\
    Ta có $\heva{&BC\perp AB\\ &BC\perp SH\ (\text{do } SH\perp(ABCD))} \Rightarrow BC\perp(SAB)$.
   \itemch Đúng. Ta có $BC\perp(SAB) \Rightarrow BC\perp HE$. \hfill(2)\\ 
   Từ $(1)$ và  $(2)$ suy ra $HE\perp(SBC)$ hay $\mathrm{d}(H,(SBC))=HE$.\\
   Tam giác $BCH$ vuông tại  $B$ có
    $$CH=\sqrt{BC^2+BH^2}=\sqrt{a^2+a^2}=a \sqrt{2}.$$
    Tam giác $SCH$ vuông tại $H$ có
   $\tan \widehat{SCH}=\dfrac{SH}{CH}  \Rightarrow SH=a \sqrt{2}$.\\
   Tam giác $SBH$ vuông tại $H$ có đường  cao $HE$ nên
   \begin{eqnarray*}
    \dfrac{1}{HE^2} &=&\dfrac{1}{SH^2}+\dfrac{1}{BH^2} \\ 
    \Rightarrow  HE&=&\dfrac{SH\cdot BH}{\sqrt{SH^2+BH^2}} \\
    &=&\dfrac{a \sqrt{2} \cdot  a}{\sqrt{2 a^2+a^2}}=\dfrac{a \sqrt{6}}{3}.
   \end{eqnarray*}
   Vậy $\mathrm{d}(H, (SBC))=HE=\dfrac{a  \sqrt{6}}{3}$.
   \itemch Đúng. Gọi $K$ là trung điểm $CD$ thì $HK$ là đường trung bình của hình chữ nhật $ABCD$  nên ${HK\parallel BC\parallel AD\Rightarrow HK\perp CD}$.\\
   Ta có ${\heva{&CD\perp HK\\  &CD\perp SH} \Rightarrow CD\perp(SHK)}$.
   \itemch  Sai. Kẻ đường cao $HI$ của tam giác $SHK$.\\
   Ta có  $\heva{&HI\perp SK\\ &HI\perp CD\ (\text{do } CD\perp(SHK), HI\subset(SHK))} \Rightarrow HI\perp(SCD)$.\\
    Tam giác $SHK$ vuông tại $H$ có đường cao $HI$ nên
   $$\dfrac{1}{HI^2}=\dfrac{1}{SH^2}+\dfrac{1}{HK^2} \Rightarrow HI=\dfrac{SH\cdot HK}{\sqrt{SH^2+HK^2}}=\dfrac{a \sqrt{2} \cdot  a}{\sqrt{2a^2+a^2}}=\dfrac{a \sqrt{6}}{3}.$$
   Vậy  $\mathrm{d}(H, (SCD))=HI=\dfrac{a \sqrt{6}}{3}$.
  \end{itemchoice}
 } 
\end{ex}


%%%%%%%%%%%%%%%% TỰ LUẬN %%%%%%%%%%%%%%%%

% nguồn: Ok_TuanLuan_Toan11_Lan2
\begin{ex}%[1H4H5-4]%[ma:9db14ac5]
	Để làm một khung lồng đèn kéo quân hình lăng trụ lục giác $ABCDEF.A'B'C'D'E'F'$. Bình gắn hai thanh tre $A_1 D_1, F_1 C_1$ song song với  mặt phẳng đáy và cắt nhau tại $O_1$ (hình bên dưới).
	\begin{center}
		\begin{tikzpicture}[scale=1.2,font=\footnotesize,line join=round,line cap=round,>=stealth]
				\coordinate (A) at  (0,0);
				 \coordinate (B) at (1,-1);
				\coordinate (C)  at (2.5,-1);
				\coordinate (D) at (3,0); ì 
				 \coordinate (E) at ($(A)+(D)-(B)$);
				\coordinate (F) at  ($(B)+(E)-(C)$);
				\foreach \i in {A,B,C,D,E,F}{\coordinate (\i') at  ($(\i)+(0,4)$);}
				\draw (A')--(B')--(C')--(D')--(E')--(F')--cycle; 
				\draw (A)--(B)--(C)--(D); 
				 \draw (A)--(A')
					(B)--(B') 
					(C)--(C') 
					(D)--(D');  (B')--(C')--(D');
				 \draw[dashed] (A)--(F)--(F')
					(D)--(E)--(E') 
					 (E)--(F);
				\foreach \i/\g in  {A'/170,B'/90,C'/90,D'/0,A/180,B/-90,C/-90,D/0,E/-90,F/-90,E'/90,F'/90}{\draw[fill=blue](\i) circle (1pt)
					($(\i)+(\g:5mm)$)  node[scale=1]{$\i$};}
				\coordinate (A1) at  ($(A)!0.5!(A')$);
				 \coordinate (D1) at ($(D)!0.5!(D')$);
				\coordinate (F1) at  ($(F)!0.5!(F')$);
				 \coordinate (C1) at ($(C)!0.5!(C')$); C
				\draw[dashed] (A1)--(D1) 
					(C1)--(F1); 
				\coordinate (Q) at  ($(A1)!0.5!(D1)$);
				\foreach  \i/\g in {A1/170,D1/0,C1/0,F1/180,Q/90}{\draw[fill=blue](\i) circle (1pt)
					($(\i)+(\g:5mm)$) node[scale=1]{$\i$};} 
			\end{tikzpicture}
	\end{center}
	\begin{enumerate}
		\item Xác định giao tuyến của mặt phẳng $\left(A_1  D_1, F_1 C_1 \right)$ với các mặt bên của lăng trụ.
		\item Cho  biết $A'A_1=6AA_1$ và $AA'=70$ cm, tính $CC_1$ và $C_1 C'$.
	\end{enumerate}
	\loigiai{ 
		\begin{center}
			\begin{tikzpicture}[scale=1.2,font=\footnotesize,line join=round,line cap=round,>=stealth]
					\coordinate (A)  at (0,0);
					\coordinate (B) at  (1,-1);
					\coordinate (C) at  (2.5,-1);
					\coordinate  (D) at (3,0);
					\coordinate (E)  at ($(A)+(D)-(B)$);
					 \coordinate (F) at ($(B)+(E)-(C)$);
					\foreach \i in {A,B,C,D,E,F}{\coordinate (\i') at  ($(\i)+(0,4)$);}
					\draw  (A')--(B')--(C')--(D')--(E')--(F')--cycle;
					\draw  (A)--(B)--(C)--(D);
					\draw (A)--(A') 
						 (B)--(B')
						(C)--(C') 
						(D)--(D'); (B')--(C')--(D'); 
					\draw[dashed] (A)--(F)--(F') 
						 (D)--(E)--(E')
						 (E)--(F);
					\foreach  \i/\g in {A'/170,B'/90,C'/90,D'/0,A/180,B/-90,C/-90,D/0,E/-90,F/-90,E'/90,F'/90}{\draw[fill=blue](\i) circle (1pt)
						($(\i)+(\g:5mm)$)  node[scale=1]{$\i$};}
					 \coordinate (A1) at ($(A)!0.5!(A')$);
					\coordinate  (D1) at ($(D)!0.5!(D')$);
					 \coordinate (F1) at ($(F)!0.5!(F')$);
					\coordinate (C1)  at ($(C)!0.5!(C')$); C
					 \draw[dashed] (A1)--(D1)
						 (C1)--(F1);
					 \coordinate (Q) at ($(A1)!0.5!(D1)$);
					 \foreach \i/\g in {A1/170,D1/0,C1/0,F1/180,Q/90}{\draw[fill=blue](\i) circle (1pt)
						 ($(\i)+(\g:5mm)$) node[scale=1]{$\i$};}
				\end{tikzpicture}
		\end{center}
		\begin{enumerate}
			\item Qua $O_1$ kẻ $B_1 C_1$ song song với mặt phẳng đáy như hình vẽ, khi  đó ta có\\
			$\left(A_1 D_1, F_1 C_1 \right)\cap \left(ABB'A'\right)=A_1  B_1$.\\
			$\left(A_1 D_1,  F_1 C_1 \right)\cap \left(BCC'B'\right)=B_1 C_1.
			 $\\
			$\left(A_1 D_1, F_1  C_1 \right)\cap \left(CDD'C'\right)=C_1 D_1.
			 $\\
			 $\left(A_1 D_1, F_1 C_1 \right)\cap \left(DEE'D'\right)=D_1 E_1.
			 $\\
			$\left(A_1  D_1, F_1 C_1 \right)\cap \left(EFF'E'\right)=E_1 F_1.
			$\\ 
			$\left(A_1  D_1, F_1 C_1 \right)\cap \left(AFF'A'\right)=A_1 F_1.
			 $
			\item Ta có $\heva{A'A_1&=6AA_1 \\  A'A_1+AA_1&=AA'=70} \Leftrightarrow \heva{A'A_1&=60 \\ AA_1&=10.}$\\
			Suy ra $CC_1=AA_1=10$  cm và $C_1 C'=A_1 A'=60$ cm.
		\end{enumerate}
	} 
\end{ex}
